\section{Conventions and preliminaries}\label{sec:setup}
\subsection{Conventions}\label{subsec:conventions}
\begin{itemize}
\item Throughout, let $k$ be the base field, of characteristic $p>0$.
\item Surfaces and curves are integral and of finite type over $k$, unless stated otherwise.
\item For a line bundle $\mathcal L$ on a proper curve over a field $F$, we write $\deg_F(\mathcal L)$ for its degree over $F$.
  For a Cartier divisor $D$ on a surface $X$ and an integral curve $C\subset X$ proper over $F$, we put $(D\cdot C)_F := \deg_F(\O_X(D)|_C)$ and write $(D^2)_F$ for self-intersection.
  Unless indicated otherwise, intersection numbers on a surface over $F$ are computed over $F$, and intersections with vertical curves on a model over a discrete valuation ring are computed over its residue field.
  We omit the subscript when the field is clear.
\item If a curve $C$ is proper over a finite extension $\ell/k$, then $\deg_k(\mathcal L) = [\ell:k]\deg_\ell(\mathcal L)$ for every line bundle $\mathcal L$ on $C$.
  For an integral proper curve $C\subset X$ over $k$, with normalization $\nu\from C^\nu\to C$ and $\ell = H^0(C^\nu,\O_{C^\nu})$, we also write $(D\cdot C)_\ell := \deg_\ell\bigl(\nu^*(\O_X(D)|_C)\bigr)$, so
  \[
    (D\cdot C)_k = [\ell:k](D\cdot C)_\ell.
  \]
\item On a normal projective surface we use Mumford's rational intersection pairing of Weil divisors, computed over the base field of the surface.
\item We write $\sim$ for linear equivalence and $\equiv$ for numerical equivalence.
\item For a $\QQ$-factorial normal projective surface $Z$, we write $N^1(Z)_\QQ$ for the $\QQ$-vector space of Weil divisors with rational coefficients modulo numerical equivalence, and $\rho(Z) = \dim N^1(Z)_\QQ$.
\item The intersection form on $N^1(Z)_\QQ$ is \emph{isotropic} if $D^2 = 0$ for some nonzero $D\in N^1(Z)_\QQ$, and \emph{anisotropic} otherwise.
\item A \emph{negative curve} on a surface is a curve with negative self-intersection.
\item A \emph{fibration} is a proper morphism $f\from W\to B$ of schemes with $f_*\O_W = \O_B$.
\item For a scheme $W$ over $k$ and a field extension $k'/k$ we write $W_{k'} = W\times_k\Spec k'$.
\item If $W$ is geometrically integral over $k$, its \emph{geometric normalization} is the normalization of $W_{\bar k}$.
\item A locally Noetherian scheme $W$ over a field $k$ is \emph{geometrically regular} at a point $w$ if, for every finite extension $k'/k$, the scheme $W_{k'}$ is regular at every point over $w$ \cite[D\'efinition~6.7.6]{EGA}.
  If $W_{\bar k}$ is locally Noetherian, for instance if $W$ is locally of finite type over $k$, then $W$ is geometrically regular if and only if $W_{\bar k}$ is regular, by \cite[Proposition~5.13.7]{EGA} and faithfully flat descent of regularity \cite[Tag~\stackstag{07NG}]{Stacks}; cf.\ \cite[Proposition~6.7.7]{EGA}.
\item We write $F^e_W$ for the $e$-th power of the absolute Frobenius morphism of an $\FF_p$-scheme $W$.
\item For a finite surjective morphism $g\from W_1\to W_2$ of normal Noetherian integral schemes and a Weil divisor $D$ on $W_2$, the pullback $g^*D$ is defined as the closure of the pullback of $D$ restricted to the regular locus of $W_2$; the complement of the regular locus has codimension at least two, and so does its preimage.
This pullback is compatible with composition, because this may be checked at points of codimension one, where all Weil divisors are Cartier.
\item For a local ring $(R,\mathfrak m,\kappa)$, we write $\widehat R$ for its $\mathfrak m$-adic completion.
\item
  For a point $z$ of a scheme $Z$, we write $\O^{\rm sh}_{Z,z}$ for the \emph{strict henselization} of $\O_{Z,z}$ with respect to a separable closure of $\kappa(z)$.
  The ring $\O^{\rm sh}_{Z,z}$ is a filtered colimit of \'etale $\O_{Z,z}$-algebras, and the map $\O_{Z,z}\to\O^{\rm sh}_{Z,z}$ is flat and local \cite[Tag~\stackstag{03QL}]{Stacks}.
\end{itemize}

\subsection{Purely inseparable base change}\label{subsec:canonical-base-change}
We record the behavior of canonical divisors under base change; cf.\ \cite{Tan18BC} and \cite[Theorems~1.1 and~1.2]{PW22}.
For our application, we need the case where the scheme $X/k$ is not necessarily of finite type, which is due to \cite{NT24}. For simplicity, we only consider the case where $X$ is geometrically integral.
\begin{lemma}\label{lem:canonical-base-change}
Let $A$ be an excellent Gorenstein Noetherian $k$-algebra of finite Krull dimension.
Let $X\to\Spec A$ be a flat Gorenstein morphism of finite type, with $X$ normal and geometrically integral over $k$.
For a finite field extension $k'/k$, put $A' = A\otimes_k k'$ and $X_{A'} = X\times_A\Spec A'$.
Let $\nu\from X'\to X_{A'}$ be the normalization of the integral scheme $X_{A'} = X_{k'}$.
Then the divisorial conductor of $\nu$ is $(p-1)C_{k'}$ for an effective integral Weil divisor $C_{k'}$, and
\begin{equation}\label{eq:canonical-base-change}
  \omega_{X'/A'}\cong\O_{X'}\bigl(-(p-1)C_{k'}\bigr)\otimes\nu^*\omega_{X_{A'}/A'},
\end{equation}
where $\omega_{X'/A'}$ denotes the relative canonical sheaf, a rank-one reflexive sheaf, and $C_{k'} = 0$ if and only if $X_{A'}$ is normal.
\end{lemma}

\begin{proof}
The normalization $\nu$ is finite since $A$ is excellent.
The ring $A'$ is Gorenstein by base change, and $X_{A'}$ is Gorenstein because $X_{A'}\to\Spec A'$ is Gorenstein \cite[Tags~\stackstag{0C03} and~\stackstag{0BJL}]{Stacks}.
Since $X$ is of finite type over the excellent ring $A$, it is Noetherian and Japanese, so \cite[Corollary~5.10]{NT24} shows that the divisorial conductor of $\nu$ is $(p-1)C_{k'}$ for an effective integral Weil divisor $C_{k'}$.
Since $\omega_{X_{A'}/A'}$ is invertible, finite duality gives \eqref{eq:canonical-base-change} in codimension one, and both sides are reflexive on $X'$, so the isomorphism extends globally.
Finally, $C_{k'} = 0$ precisely when $X_{A'}$ is normal in codimension one; since $X_{A'}$ is Gorenstein and hence $S_2$, this is equivalent to normality by Serre's criterion.
\end{proof}

For extensions of height one, the vanishing of this conductor divisor admits the following local criterion, which can be seen as an algebraic version of Ji--Waldron's description of normalized base changes by foliations \cite[Sections~3 and~4]{JW21}.
\begin{lemma}\label{lem:normality-derivation}
Let $A$ be a normal $k$-domain, and let $k'/k$ be finite with $k'^p\subseteq k$.
Suppose that $B := A\otimes_k k'$ is a domain, with normalization $\widetilde B$.
For $D\in\Der_k(k')$, write $D_B = \mathrm{id}_A\otimes D$ and use the same notation for its extension to $\Frac(B)$.
Then $B$ is normal if and only if $D_B$ preserves $\widetilde B$ for every $D\in\Der_k(k')$.
\end{lemma}

\begin{proof}
The forward implication is immediate.
For the converse, put $L = \Frac(A)$ and $K = L\otimes_k k' = \Frac(B)$.
Choose a $p$-basis $\gamma_1,\ldots,\gamma_r$ of $k'/k$ and let $\delta_i$ be the $L$-derivation of $K$ with $\delta_i(\gamma_j) = \delta_{ij}$.
The commuting derivations $\delta_i$ satisfy $\delta_i^p = 0$ and preserve $\widetilde B$ by hypothesis.
They therefore define an action of $G = \alpha_p^r$ on $\Spec\widetilde B$ \cite[Section~4]{ST23}.
Writing $t_1,\ldots,t_r$ for the coordinates on $G$, its coaction on $k'$ sends $\gamma_i$ to $\gamma_i+t_i$, so the inclusion $k'\hookrightarrow\widetilde B$ is equivariant.
Since $\Spec k'\to\Spec k$ is a $G$-torsor, faithfully flat descent for algebras gives
\[
  \widetilde B\cong (\widetilde B)^G\otimes_k k'.
\]
The common kernel of the $\delta_i$ in $K$ is $L$, so $(\widetilde B)^G = \widetilde B\cap L$.
Every element of this intersection is integral over $A$, since $B$ is integral over $A$.
Normality of $A$ therefore gives $(\widetilde B)^G = A$, and the displayed isomorphism is the natural map $B\to\widetilde B$.
Hence $B$ is normal.
\end{proof}


\subsection{Killing a cohomology class by a generic $\alpha_p$-torsor}\label{subsec:generic-alpha-p-torsor}
\begin{lemma}\label{lem:generic-alpha-p-torsor}
Let $X/k$ be a regular projective geometrically integral variety with $H^1(X,\O_X)\neq 0$.
Suppose that $H^1((X_{k^{1/p^e}})^\nu,\O) = 0$ for some $e\geq 1$.
Then over $K = k(a)$, with $a$ transcendental over $k$, there is an $\alpha_p$-torsor $\pi\from V\to X_K$, where $V$ is a regular projective geometrically non-reduced variety with $H^0(V,\O_V) = K$ and $\omega_V\cong\pi^*\omega_{X_K}$.
\end{lemma}

\begin{proof}
Put $\ell = k^{1/p^e}$ and $T = (X_\ell)^\nu$.
Then $H^0(T,\O_T) = \ell$.
The absolute Frobenius $F_X^e$ factors through $T\to X$.
Since $H^1(T,\O_T) = 0$, we have $(F_X^*)^e = 0$ on $H^1(X,\O_X)$.
Choose $\eta\neq 0$ in $H^1(X,\O_X)$ and let $r\geq 1$ be minimal with $(F_X^*)^r\eta = 0$.
Then $\xi := (F_X^*)^{r-1}\eta$ is nonzero and satisfies $F_X^*(\xi) = 0$.
Choose a finite affine open cover $\{U_i = \Spec A_i\}$ and a \v{C}ech cocycle $(g_{ij})$ representing $\xi$.
Since $F_X^*(\xi) = 0$, the cocycle $(g_{ij}^p)$ is a coboundary, so there are $f_i\in A_i$ with $g_{ij}^p = f_i-f_j$.
Write $\xi_K = \xi\otimes 1\in H^1(X_K,\O_{X_K})$.

Glue
\[
  V_i = \Spec\bigl((A_i\otimes_kK)[t_i]/(t_i^p-f_i-a)\bigr)
\]
by $t_i = t_j+g_{ij}$.
These transitions preserve the equations because $g_{ij}^p = f_i-f_j$, and they satisfy the cocycle condition.
Translations $t_i\mapsto t_i+\varepsilon$ define an $\alpha_p$-torsor $\pi\from V\to X_K$.
Since $g_{ij} = t_i-t_j$ on $V$, we have $\pi^*\xi_K = 0$ in $H^1(V,\O_V)$.
Each $V_i$ is regular, since its coordinate ring is a localization of the regular ring
\[
  A_i[a,t_i]/(t_i^p-f_i-a)\cong A_i[t_i].
\]
Hence $V$ is regular.
Since $\pi$ is a universal homeomorphism and $X_K$ is integral, $V$ is integral.
The morphism $\pi$ is finite, so $V$ is projective over $K$.

To show that $H^0(V,\O_V) = K$, let $s\in H^0(V,\O_V)$ be nonzero, and represent $s|_{V_i}$ by a polynomial $s_i(t_i)\in (A_i\otimes_kK)[t_i]$ of degree $<p$.
The transitions preserve degree, so let $m$ be the common degree of the $s_i$.
Their leading coefficients glue to an element $b\in H^0(X_K,\O_{X_K}) = K$.
If $m>0$, write their two leading terms as $bt_i^m+c_it_i^{m-1}$, where $b\in K^\times$ and $c_i\in A_i\otimes_kK$.
Comparing coefficients of $t_j^{m-1}$ in $s_i(t_j+g_{ij}) = s_j(t_j)$ gives $c_j-c_i = mbg_{ij}$, so $mb\xi_K = 0$ in $H^1(X_K,\O_{X_K})$.
This contradicts $\xi_K\neq 0$.
Thus the $s_i$ have degree zero and glue to an element of $K$, so $H^0(V,\O_V) = K$.

Let $W_i$ be the inverse image of $U_i$ under $T\to X$.
Since $H^1(T,\O_T) = 0$, the pulled-back cocycle has the form $g_{ij} = h_i-h_j$ with $h_i\in H^0(W_i,\O_T)$.
The functions $f_i-h_i^p$ glue to a constant $\lambda\in H^0(T,\mathcal O_T) = \ell$.
Put $L = \ell(a)$ and $T_L = T\times_\ell L$.
Consider the Cartesian square
\[
  \begin{array}{ccc}
    V\times_{X_K}T_L &\longrightarrow& V\\
    \downarrow &&\downarrow\pi\\
    T_L &\longrightarrow& X_K.
  \end{array}
\]
The coordinates $z_i = t_i-h_i$ give an isomorphism
\[
  V\times_{X_K}T_L\cong T_L\times_L\Spec L[z]/(z^p-a-\lambda).
\]
Put $L' = L((a+\lambda)^{1/p})$ and $T_{L'} = T\times_\ell L'$.
Then the projection $V\times_{X_K}T_{L'}\to T_{L'}$ is a trivial $\alpha_p$-torsor.
Since $T_{L'}\to X_{L'}$ is birational, $V_{L'}$ is generically nonreduced.
Hence $V$ is not geometrically reduced over $K$.

Put $B_i = H^0(V_i,\O_{V_i})$.
Let $\varphi_i\from B_i\to A_i\otimes_kK$ be the coefficient functional $\sum_{r<p}b_rt_i^r\mapsto b_{p-1}$.
For $0\leq r<p$, the multiple $t_i^{p-1-r}\varphi_i$ extracts the coefficient of $t_i^r$.
These multiples form the dual basis to $1,t_i,\ldots,t_i^{p-1}$, so $\varphi_i$ generates $\operatorname{Hom}_{A_i\otimes_kK}(B_i,A_i\otimes_kK)$ as a $B_i$-module.
The translations $t_i = t_j+g_{ij}$ preserve the coefficient of degree $p-1$, so these generators glue to trivialize $\omega_{V/X_K}$.
Finite duality then gives $\omega_V\cong\pi^*\omega_{X_K}$.
\end{proof}

\subsection{Tangent cones}\label{subsec:tangent-cones}
Let $X$ be a variety over $k$, let $x\in X$ be a closed point at which $X$ is regular, with $\kappa(x)$ purely inseparable over $k$, and let $\bar x\in X_{\bar k}$ be the geometric point over $x$.
Let $\mathfrak m_{\bar x}$ be the maximal ideal of $\O_{X_{\bar k},\bar x}$, and let $T = (\mathfrak m_{\bar x}/\mathfrak m_{\bar x}^2)^\vee$ be the Zariski tangent space of $X_{\bar k}$ at $\bar x$.
Write $C = \Spec\bigl(\bigoplus_{j\geq 0}\mathfrak m_{\bar x}^j/\mathfrak m_{\bar x}^{j+1}\bigr)\subseteq T$ for the tangent cone and $I_C$ for its homogeneous ideal.
The tangent cone $C$ has the same dimension as $X$, and $C = T$ if $\bar x$ is regular.
If $\bar x$ is a singular point, let $e$ be the lowest degree of a nonzero form in $I_C$.
We define the space of \emph{vertical directions} to be
\[
  L = \{v\in T : v(f) = 0\text{ for all }f\in\mathfrak m_x\},
\]
where $\mathfrak m_x$ is the maximal ideal of $\O_{X,x}$, mapped to $\mathfrak m_{\bar x}/\mathfrak m_{\bar x}^2$ by the projection $\pi\from X_{\bar k}\to X$.
Equivalently, $L$ is the Zariski tangent space at $\bar x$ of the one-point fibre $\pi^{-1}(x) = \Spec(\kappa(x)\otimes_k\overline k)$.

\begin{lemma}\label{lem:tangent-cone}
In this situation, the following hold.
\begin{enumerate}[label=\textup{(\arabic*)}]
\item $\dim L = \dim_{\kappa(x)}\Omega_{\kappa(x)/k}$ and $\dim T\leq\dim X+\dim L$; in particular, $\dim T\leq\dim X+1$ if $\kappa(x)/k$ is simple.
\item If $y_1,\dots,y_N$ is a basis of $\mathfrak m_{\bar x}/\mathfrak m_{\bar x}^2$ such that $y_1,\dots,y_u$ span the image of $\mathfrak m_x$, so that $L = \{v\in T : v(y_1) = \dots = v(y_u) = 0\}$, then $I_C$ is stable under $\partial/\partial y_j$ for $j>u$.
\item Suppose that $e<p$, and let $y_1,\dots,y_N$ be as in \textup{(2)}.
  Then every form of degree $e$ in $I_C$ is a polynomial in $y_1,\dots,y_u$ alone, so it does not involve the coordinates along $L$.
  In particular, if $I_C$ is generated by one form, then $C$ is the preimage of a cone in $T/L$ under the projection $T\to T/L$.
\end{enumerate}
\end{lemma}

\begin{proof}
Put $\kappa = \kappa(x)$.
For (1), the description of $L$ as the tangent space of the fibre gives $L^\vee\cong\Omega_{\kappa/k}\otimes_\kappa\bar k$.
Since $\mathfrak m_x$ is generated by $\dim X$ elements and $L$ is the annihilator of its image in the cotangent space, $\dim T-\dim L\leq\dim X$.

For (2), put $A = \O_{X,x}$ and $B = A\otimes_k\kappa$.
The ring $B$ is finite over $A$, and $B/\mathfrak m_xB = \kappa\otimes_k\kappa$ is local with residue field $\kappa$ because $\kappa/k$ is purely inseparable.
Hence $B$ is local with residue field $\kappa$, and it is the local ring of $X_\kappa$ at the unique point over $x$, which is $\kappa$-rational.
Let $\mathfrak m$ be the maximal ideal of $B$, let $\Sigma\subseteq\mathfrak m/\mathfrak m^2$ be the $\kappa$-span of the image of $\mathfrak m_x$, and write the associated graded ring $\operatorname{gr}_{\mathfrak m}B$ as $S/I^*$ with $S = \Sym_\kappa(\mathfrak m/\mathfrak m^2)$, where $I^*$ is the ideal of initial forms.
The maximal ideal of $\O_{X_{\bar k},\bar x}$ is generated by $\mathfrak m$, and $\operatorname{gr}_{\mathfrak m_{\bar x}}\O_{X_{\bar k},\bar x} = \operatorname{gr}_{\mathfrak m}B\otimes_\kappa\overline k$.
Hence $T = (\mathfrak m/\mathfrak m^2)^\vee\otimes_\kappa\overline k$, $I_C = I^*\otimes_\kappa\overline k$, and $L$ is the annihilator of $\Sigma$.

For $b\in\kappa$ with a lift $\tilde b\in A$, put $w_b = \tilde b\otimes 1-1\otimes b$, which lies in $\mathfrak m$.
Let $\delta\in\Der_k(\kappa)$.
Then $D = \mathrm{id}_A\otimes\delta$ is a derivation of $B$ which vanishes on $\mathfrak m_x$ and satisfies $D(w_b) = -1\otimes\delta(b)$, so $D(\mathfrak m^j)\subseteq\mathfrak m^{j-1}$ for $j\geq 1$.
The induced maps $\mathfrak m^j/\mathfrak m^{j+1}\to\mathfrak m^{j-1}/\mathfrak m^j$ form a derivation of degree $-1$ of $S/I^*$, which is $\kappa$-linear because $f\,D(1\otimes c)\in\mathfrak m^j$ for $f\in\mathfrak m^j$ and $c\in\kappa$.
Hence it is induced by the $\kappa$-linear derivation $\partial_\lambda$ of $S$ with $\partial_\lambda(y) = \lambda(y)$ for $y\in\mathfrak m/\mathfrak m^2$, where $\lambda$ is the functional $f\mapsto D(f)\bmod\mathfrak m$ on $\mathfrak m/\mathfrak m^2$, and therefore $\partial_\lambda(I^*)\subseteq I^*$.
The functional $\lambda$ vanishes on $\Sigma$, and $\lambda(w_b) = -\delta(b)$.

The cotangent-space identification for the fibre $\Spec(\kappa\otimes_k\kappa)$ gives an isomorphism $\Omega_{\kappa/k}\cong(\mathfrak m/\mathfrak m^2)/\Sigma$, sending $db$ to the class of $w_b$.
Under this identification, $\lambda$ corresponds to $-\delta$, so these functionals exhaust the dual of $(\mathfrak m/\mathfrak m^2)/\Sigma$ as $\delta$ varies.
Thus $I^*$ is stable under $\partial_\lambda$ for every functional $\lambda$ on $\mathfrak m/\mathfrak m^2$ vanishing on $\Sigma$.

Choose a basis $y_1,\dots,y_N$ of $\mathfrak m/\mathfrak m^2$ such that $y_1,\dots,y_u$ span $\Sigma$.
Then $L = \{v\in T : v(y_1) = \dots = v(y_u) = 0\}$, and $I^*$ is stable under $\partial/\partial y_j$ for $j>u$.
The derivations $\partial/\partial y_j$ with $j>u$ are the derivatives along a basis of $L$, so this condition does not depend on the choice of coordinates.

For (3), let $Q\in I^*$ be homogeneous of degree $e<p$.
For $j>u$ we have $\partial Q/\partial y_j\in I^*_{e-1} = 0$, and every exponent occurring in $Q$ is at most $e<p$, so $Q$ is a polynomial in $y_1,\dots,y_u$.
The forms of degree $e$ in $I_C$ are the $\overline k$-linear combinations of such $Q$, and $\overline k[y_1,\dots,y_u]$ is the coordinate ring of $T/L$ for every choice of coordinates as in (2), which proves the first assertion of (3).
If $I_C$ is generated by one form, then this form has degree $e$, so $C$ is the preimage of the cone in $T/L$ defined by the same form.
\end{proof}


\subsection{Birational contractions}\label{subsec:birational-surfaces}
% Let $\varepsilon\from Y'\to Y$ be the blowup of a codimension-two closed point $P$ on a regular Noetherian integral scheme $Y$ of dimension two, with exceptional curve $E$.
% Then $Y'$ is regular, $E\cong\PP^1_{\kappa(P)}$ and $\O_E(E)\cong\O_{\PP^1_{\kappa(P)}}(-1)$ \cite[Tags~\stackstag{0AGQ} and~\stackstag{0AGR}]{Stacks}.
% If $Y$ admits a dualizing sheaf $\omega_Y$, adjunction gives the induced dualizing sheaf on $Y'$ as
% \begin{equation}\label{eq:blowup-canonical}
%   \omega_{Y'}\cong\varepsilon^*\omega_Y\otimes\O_{Y'}(E).
% \end{equation}
% The same formula holds for the relative dualizing sheaves of regular flat models of curves over a discrete valuation ring.
% If $Y$ is a projective surface over $k$, then, with compatible canonical divisors,
% \begin{equation}\label{eq:blowup-intersections}
%   K_{Y'} = \varepsilon^*K_Y+E,\qquad
%   (K_{Y'}\cdot E)_k = (E^2)_k = -[\kappa(P):k],\qquad
%   (K_{Y'}^2)_k = (K_Y^2)_k-[\kappa(P):k].
% \end{equation}
% Moreover, if $-K_{Y'}$ is nef, then $-K_Y$ is nef: for a curve $C\subset Y$ with strict transform $\widetilde C$, the projection formula gives
% \[
%   -K_Y\cdot C = (-K_{Y'}+E)\cdot\widetilde C\geq 0.
% \]

% Let $\mu\from T\to S$ be the minimal resolution of a normal projective surface over an algebraically closed field.
% Choose canonical divisors with $\mu_*K_T = K_S$.
% Minimality makes $K_T$ relatively nef by \cite[Theorem~10.5]{Kol13}, and the negative definite exceptional intersection matrix gives
% \begin{equation}\label{eq:minimal-resolution-canonical}
%   K_T = \mu^*K_S-Z,\qquad Z\geq 0,
% \end{equation}
% with $Z$ an exceptional $\QQ$-divisor \cite[Claim~10.3.5]{Kol13}.
% Here $\mu^*K_S$ is Mumford's numerical pullback, characterized by $\mu_*(\mu^*K_S) = K_S$ and $(\mu^*K_S)\cdot E = 0$ for every exceptional curve $E$.
% If $K_S$ is Cartier, then this is the ordinary pullback and $Z$ is integral.

\begin{definition}\label{def:minus-one-curve}
Let $X$ be a regular integral scheme of dimension two, proper and flat over $S$, where $S$ is either the spectrum of a field or a Dedekind scheme of dimension one.
A \emph{vertical curve} on $X$ is an integral closed subscheme $C\subset X$ of dimension one contained in a fibre $X_s$ over a closed point $s\in S$.
We call such a curve $C$ a \emph{$(-1)$-curve} if
\[
  (K_{X/S}\cdot C)_{k_C} = (C^2)_{k_C} = -1,
\]
where $K_{X/S}$ is a relative canonical divisor and $k_C = H^0(C,\O_C)$.
For such a curve $C$, adjunction gives $H^1(C,\O_C) = 0$.
By Castelnuovo's criterion \cite[Theorem~9.3.8 and Proposition~9.3.10]{Liu02}, these are exactly the vertical curves that can be contracted to a regular point.
\end{definition}

We record standard facts about contracting negative curves, including the residue-field details needed below.

\begin{lemma}\label{lem:contraction}
Let $k$ be a separably closed field of characteristic $p>3$, and let $X$ be a regular projective surface over $k$ with $K_X\equiv 0$.
Let $C\subset X$ be a negative curve, and put $\ell = H^0(C,\O_C)$.
\begin{enumerate}
\item $C\cong\PP^1_\ell$, $(C^2)_\ell = -2$, and $\O_C(C)\cong\O_{\PP^1_\ell}(-2)$.
\item There is a birational morphism $f\from X\to X'$ onto a normal projective surface which contracts exactly $C$, to a point $P$ with $\kappa(P) = \ell$.
  Moreover $R^1f_*\O_X = 0$ and $\rho(X') = \rho(X)-1$.
\item The surface $X'$ is Gorenstein, and $f^*\omega_{X'}\cong\omega_X$.
\item $2D$ is Cartier for every Weil divisor $D$ on $X'$.
\item There is an isomorphism of rings $\widehat\O_{X',P}\cong \ell[[u,v,w]]/(uv-w^2)$; in particular $X'$ is a local complete intersection over $k$.
  \textup{(}This isomorphism need not be an isomorphism of $k$-algebras.\textup{)}
\end{enumerate}
\end{lemma}

\begin{proof}
(1) The curve $C$ is a Cartier divisor on the regular surface $X$, hence Gorenstein.
Put $g = \dim_\ell H^1(C,\O_C)$.
Adjunction gives
\[
  (C^2)_\ell = ((K_X+C)\cdot C)_\ell = 2(g-1)<0,
\]
so $g = 0$ and $(C^2)_\ell = -2$.
By \cite[Lemma~10.6(3)]{Kol13}, $C$ is a conic or $\PP^1$ over $\ell$.
Since $\ell$ is separably closed and $p\neq 2$, an integral conic over $\ell$ is smooth and split, so $C\cong\PP^1_\ell$.
The normal bundle has degree $(C^2)_\ell = -2$, hence is $\O_{\PP^1_\ell}(-2)$, proving (1).

(2) The negative curve $C$ spans a $(K_X+C)$-negative extremal ray.
Applying \cite[Theorem~4.4]{Tan18} with boundary $C$ gives a projective birational morphism $f\from X\to X'$ contracting exactly $C$, with $X'$ normal and projective and $\rho(X') = \rho(X)-1$.
Since $K_X\cdot C = 0$ and $-C$ is $f$-nef, \cite[Theorem~10.4]{Kol13} gives $R^1f_*\O_X(-jC) = 0$ for every $j\geq 0$.
Thus $P$ is a rational singularity, and \cite[Proposition~10.9(4)]{Kol13} gives $\kappa(P) = \ell$.

(3) The numerical discrepancy of $C$ is zero because $K_X\cdot C = 0$.
Hence $X'$ is Gorenstein and $f$ is crepant by \cite[Theorem~2.29(2)]{Kol13}.

(4) The intersection matrix over $\ell$ is $(-2)$, so \cite[Proposition~10.9(3)]{Kol13} shows that $2D$ is Cartier at $P$ for every Weil divisor $D$ on $X'$.
Since $X'$ is regular away from $P$, the divisor $2D$ is Cartier everywhere.

(5) Put $R = \O_{X',P}$ and $I_j = (f_*\O_X(-jC))_P$.
Pushing forward $0\to\O_X(-C)\to\O_X\to\O_C\to0$ and taking the stalk at $P$ gives $I_1 = \ker(R\to H^0(C,\O_C)=\kappa(P)) = \mathfrak m_P$.
Similarly, since $R^1f_*\O_X(-jC) = 0$ for all $j\geq 0$, we have $I_j/I_{j+1}\cong H^0(\PP^1_\ell,\O(2j))$, compatibly with multiplication.
The multiplication map $H^0(\PP^1_\ell,\O(2))^{\otimes_\ell j}\to H^0(\PP^1_\ell,\O(2j))$ is surjective, since every monomial of degree $2j$ is a product of $j$ quadratic monomials.
Hence every class in $I_j/I_{j+1}$ is a sum of products of $j$ classes from $I_1/I_2$, so $I_j = I_1^j+I_{j+1} = \mathfrak m_P^j+I_{j+1}$.
Iterating this equality gives $I_j = \mathfrak m_P^j+I_n$ for all $n\geq j$.
By formal functions, $\widehat R\cong\varprojlim_n R/I_n$; taking limits gives $I_j\widehat R = \mathfrak m_P^j\widehat R$ and hence
\[
  \operatorname{gr}_{\mathfrak m_P}\widehat R
  \cong \bigoplus_{j\geq 0}H^0(\PP^1_\ell,\O(2j))
  \cong \ell[u,v,w]/(uv-w^2).
\]
The displayed graded ring gives $\length_R(R/\mathfrak m_P^n) = \sum_{j=0}^{n-1}(2j+1) = n^2$ for $n\geq 1$.
After choosing a coefficient field identified with $\ell$, \cite[Lemma~23.4]{Lip69} therefore gives a quadratic hypersurface presentation $\widehat R\cong\ell[[u,v,w]]/(g)$.
Comparing associated graded rings, we may choose the variables $u,v,w$ and rescale $g$ so that $g-(uv-w^2)\in(u,v,w)^3$.
Since $p\neq 2$, the quadratic form $uv-w^2$ is nondegenerate.
The formal Morse lemma \cite[Lemma~10.43]{Kol23}, applied with coefficient ring $\ell$ and maximal ideal $(0)$, therefore gives $\widehat R\cong\ell[[u,v,w]]/(uv-w^2)$ as rings.
The hypersurface presentation and regularity of $X'$ away from $P$ imply that $X'$ is a local complete intersection over $k$ \cite[Tag~\stackstag{09Q6}]{Stacks}.
\end{proof}


\subsection{Even lattices}\label{subsec:even-lattices}
An \emph{integral lattice} is a finitely generated free abelian group $N$ equipped with a nondegenerate symmetric bilinear pairing $(\ ,\ )\from N\times N\to\ZZ$.
It is \emph{even} if $(v,v)\in 2\ZZ$ for every $v\in N$.
For a $\ZZ$-basis $e_1,\ldots,e_r$ of $N$, the matrix $G = ((e_i,e_j))_{1\leq i,j\leq r}$ is called a \emph{Gram matrix} of $N$.
Its determinant $\det N := \det G$, called the \emph{Gram determinant}, is independent of the chosen $\ZZ$-basis.
The lattice $N$ is \emph{unimodular} if $\lvert\det N\rvert = 1$, equivalently if the pairing induces an isomorphism $N\cong\Hom_\ZZ(N,\ZZ)$.
For a lattice over $\ZZ_\ell$, unimodularity means that its Gram determinant is a unit in $\ZZ_\ell$.
The pairing induces a quadratic form $q(v) = (v,v)$ on $N\otimes\QQ$.
An nonzero element $v\in N\otimes\QQ$ is called \emph{anisotropic} if $q(v)\neq 0$, and \emph{isotropic} if $q(v) = 0$.
The quadratic form $q$ (or the corresponding pairing) is called \emph{anisotropic} if it has no isotropic elements, and \emph{isotropic} otherwise. %ref: ~/math/students/Merkurev/QuadForm



\begin{lemma}\label{lem:lattice-dichotomy}
Let $p$ be an odd prime and let $N$ be an even integral lattice of signature $(1,r-1)$ with $\lvert\det N\rvert = p^a$ for an integer $a\geq 0$.
Then $r$ is even, and $N\otimes\QQ$ is anisotropic if and only if $r = 2$ and $a$ is odd.
In this case $p\equiv 1\pmod 4$.
\end{lemma}

\begin{proof}
Since $N$ is even of odd determinant, the induced bilinear form on $N/2N$ with values in $\FF_2$ is nondegenerate and alternating, so $r$ is even.
In particularly, $r\ge2$.
Then the quadratic form on $N\otimes\QQ$ is indefinite by the signature.
If $r\geq 5$, then this implies that the form is isotropic by Meyer's theorem \cite[Chapter~IV, Section~3.2, Corollary~2 to Theorem~8]{Ser73}.

Suppose that $r = 4$.
By the Hasse--Minkowski theorem \cite[Chapter~IV, Section~3.2, Theorem~8]{Ser73}, to show that $N\otimes\QQ$ is isotropic, it suffices to show that $N\otimes\RR$ is isotropic and that $N\otimes\QQ_\ell$ is isotropic for every prime $\ell$.

At $\ell = 2$, $\det N$ is a unit in $\ZZ_2$ because $\lvert\det N\rvert = p^a$ and $p$ is odd, so $N\otimes\ZZ_2$ is \emph{unimodular}.
Since it is also even, it is an orthogonal sum of two binary lattices with Gram matrices \cite[Hilfssatz~5]{Cas62}
\[
  B_1 = \begin{pmatrix}2a_1&1\\1&2b_1\end{pmatrix},\qquad
  B_2 = \begin{pmatrix}2a_2&1\\1&2b_2\end{pmatrix},
\]
where $a_1,b_1,a_2,b_2\in\ZZ_2$.
Since $\det B_1\equiv\det B_2\equiv -1\pmod 4$ and squares of $2$-adic units are $1$ modulo $4$, we have $\det N\equiv 1\pmod 4$.
Since $\det N = -p^a$, we get $p\equiv -1\pmod 4$ and that $a$ is odd.
In these coordinates, the quadratic form satisfies $q = 2Q$, where
\[
  Q(x_1,x_2,x_3,x_4) = a_1x_1^2+x_1x_2+b_1x_2^2+a_2x_3^2+x_3x_4+b_2x_4^2.
\]
The polynomial $Q$ has coefficients in $\ZZ_2$ and has the same zeros as $q$ over $\QQ_2$.
The values of $Q$ at $(1,0,0,0)$, $(0,0,1,0)$ and $(1,0,1,0)$ are $a_1$, $a_2$ and $a_1+a_2$, respectively.
So modulo $2$, the form $Q$ has a zero at one of these points.
The gradient of $Q$ modulo $2$ is $(x_2,x_1,x_4,x_3)$, so this zero is nonsingular and lifts by Hensel's lemma to a nonzero isotropic vector over $\QQ_2$.

At the real place, $N\otimes\RR$ has signature $(1,3)$ and is therefore isotropic.
At an odd prime $\ell\neq p$, the determinant is a unit in $\ZZ_\ell$, so the reduction of the lattice is a nondegenerate quaternary form over $\FF_\ell$.
Since $\ell$ is odd, \cite[Chapter~IV, Section~1.4, Theorem~1]{Ser73} gives coordinates in which this reduction is $\bar q = \sum_{i=1}^4 c_ix_i^2$ with $c_i\in\FF_\ell^\times$.
The sets $\{c_1x^2:x\in\FF_\ell\}$ and $\{-c_3-c_2y^2:y\in\FF_\ell\}$ each have $(\ell+1)/2$ elements, so they intersect.
Thus $\bar q(x,y,1,0) = 0$ for some $x,y\in\FF_\ell$, and the partial derivative with respect to $x_3$ at this zero is $2c_3\neq 0$.
Hensel's lemma therefore lifts this zero to a nonzero isotropic vector over $\QQ_\ell$.
At $\ell = p$, an anisotropic quaternary quadratic form has square discriminant \cite[Chapter~IV, Section~2.2, Theorem~6(iii)]{Ser73}, whereas $\det N$ has $p$-adic valuation $a$, which is odd, a contradiction.

Suppose $r = 2$.
Choose an orthogonal basis of $N\otimes\QQ$ and write $q(x,y) = ux^2+vy^2$ with $u,v\in\QQ^\times$.
A nonzero zero must have $y\neq 0$, and $q(x,y) = 0$ is then equivalent to $(x/y)^2 = -v/u$.
Thus $N\otimes\QQ$ is isotropic if and only if $-v/u$ is a rational square, or equivalently if $-uv$ is a rational square.
A change of basis multiplies the Gram determinant by a rational square, so $uv$ and $\det N$ have the same square class, and the criterion is equivalent to $-\det N$ being a rational square.
The signature $(1,1)$ gives $\det N<0$, so $-\det N = p^a$ is a rational square if and only if $a$ is even.
Therefore, under the assumption that $N \otimes \mathbb Q$ is anisotropic, $a$ is odd.
Since $N$ is even, its Gram matrix in an integral basis has the form $\begin{pmatrix}2b&c\\c&2d\end{pmatrix}$ with $b,c,d\in\ZZ$.
Its determinant $4bd-c^2$ is odd, so $c$ is odd and $\det N\equiv -1\pmod 4$.
Since $\det N = -p^a$, this gives $p^a\equiv 1\pmod 4$.
As $a$ is odd, $p^a\equiv p\pmod 4$, hence $p\equiv 1\pmod 4$.
\end{proof}


\subsection{Genus-one fibrations}
\begin{lemma}\label{lem:fibration}
Let $X$ be a regular projective surface over $k$ with $H^0(X,\O_X) = k$, $K_X\equiv 0$ and $\chi(\O_X) = 2$.
If the intersection form on $N^1(X)_\QQ$ is isotropic, then there is a fibration $f\from X\to B$ onto a regular projective curve.
\end{lemma}

\begin{proof}
Choose an ample divisor $A$ and a divisor $D$ with $D^2 = 0$ and $A\cdot D>0$.
Riemann--Roch gives $h^0(D)+h^0(K_X-D)\geq\chi(\O_X(D)) = \chi(\mathcal O_X) = 2$.
Hence $h^0(D)\geq 2$.

Suppose that $D$ is not nef, and let $C$ be a curve with $D\cdot C<0$.
Then $C$ is a component of every member of $\lvert D\rvert$, so $C^2<0$.
As in the proof of Lemma~\ref{lem:contraction}(1), with $h = [H^0(C,\O_C):k]$ we have $(C^2)_k = -2h$, and $(D\cdot C)_k\in h\ZZ$.
Hence $r_C(D) := D+\frac{(D\cdot C)_k}{h}C$ is a $\ZZ$-divisor with $(r_C(D)^2)_k = 0$.
The reflection $r_C$ is an isometry of $N^1(X)_\QQ$ preserving the component of the positive cone containing $A$; hence $A\cdot r_C(D)>0$, while $A\cdot r_C(D)<A\cdot D$.
By Riemann--Roch again, $h^0(r_C(D))\geq 2$.
Since the positive integer $A\cdot D$ decreases, after finitely many reflections we obtain a nef divisor $D$ with $D^2 = 0$ and $h^0(D)\geq 2$.

It follows that $\kappa(X,D) = 1$, so the nef divisor $D$ is semiample (cf.\ \cite[Lemma~2.7]{Tan20}).
The Stein factorization of the morphism defined by a sufficiently divisible multiple of $D$ gives a fibration $f\from X\to B$ onto a regular projective curve.
\end{proof}
